\documentclass{article}
\usepackage{fancyhdr}
\usepackage{import}
\usepackage{tcolorbox}
\usepackage{amsmath}
\usepackage{amssymb}
\usepackage{amsthm}
\usepackage{geometry}
\usepackage{enumitem}
\usepackage{pdfpages}
\usepackage{transparent}
\usepackage[hidelinks]{hyperref}
\usepackage{booktabs}
\usepackage{bm}
\usepackage{tabularx}
\tcbuselibrary{theorems}
\tcbuselibrary{skins}

% tabularx Settings
\renewcommand\tabularxcolumn[1]{>{\centering\arraybackslash}m{#1}}

% Page Settings
\newgeometry{vmargin={15mm}, hmargin={15mm, 15mm}}
\setlength{\parindent}{0pt}
\setlength{\parskip}{1em plus 1pt minus 1pt}

% hyperref Settings
\hypersetup{
    colorlinks=true,
    linkcolor=blue,
    filecolor=magenta,      
    urlcolor=cyan,
}

% inkscape Settings
\newcommand{\incfig}[1]{%
    \def\svgwidth{\columnwidth}
    \import{./figure/}{#1.pdf_tex}
}

% Definition Box Styles
\newtcbtheorem[auto counter, number within=subsection]{defn}{Definition}{
% basic settings
    theorem name and number,
    separator sign={},
    description delimiters={(}{)},
    enhanced jigsaw,
    sharp corners,
% box margins and rules
    boxrule=1pt,
    left=0.2cm,right=0.2cm,top=0.2cm,
    toptitle=0.1cm,
    bottomtitle=0.08cm,
% box colors
    colframe=white!25!black,
    colback=white,
    coltitle=white,
% box fonts
    fonttitle=\bfseries,fontupper=\normalsize
}{defn}

% Theorem Box Styles
\newtcbtheorem[auto counter, number within=subsection]{thm}{Theorem}{
% basic settings
    theorem name and number,
    separator sign={},
    description delimiters={(}{)},
    enhanced jigsaw,
    sharp corners,
% box margins and rules
    boxrule=1pt,
    left=0.2cm,right=0.2cm,top=0.2cm,
    toptitle=0.1cm,
    bottomtitle=0.08cm,
% box colors
    colframe=white!25!black,
    colback=white,
    coltitle=white,
% box fonts
    fonttitle=\bfseries,fontupper=\normalsize
}{thm}

% Lemma Box Styles
\newtcbtheorem[number within=tcb@cnt@thm]{lem}{Lemma}{
% basic settings
    theorem name and number,
    separator sign={},
    description delimiters={(}{)},
    enhanced jigsaw,
    sharp corners,
% box margins and rules
    boxrule=1pt,
    left=0.2cm,right=0.2cm,top=0.2cm,
    toptitle=0.1cm,
    bottomtitle=0.08cm,
% box colors
    colframe=white!50!black,
    colback=white,
    coltitle=white,
% box fonts
    fonttitle=\bfseries,fontupper=\normalsize
}{lem}

% Corollary Box Styles
\newtcbtheorem[number within=tcb@cnt@thm]{cor}{Corollary}{
% basic settings    
    theorem name and number,
    separator sign={},
    description delimiters={(}{)},
    enhanced jigsaw,
    sharp corners,
% box margins and rules
    boxrule=1pt,
    left=0.2cm,right=0.2cm,top=0.2cm,
    toptitle=0.1cm,
    bottomtitle=0.08cm,
% box colors
    colframe=white!50!black,
    colback=white,
    coltitle=white,
% box fonts
    fonttitle=\bfseries,fontupper=\normalsize
}{cor}

% Remark Box Styles
\newtcbtheorem[number within=subsection]{rem}{Remark}{
% basic settings
    theorem name and number,
    separator sign={},
    description delimiters={(}{)},
    enhanced jigsaw,
    sharp corners,
% box colors
    colframe=white!60!black,
    colback=white,
    coltitle=black,
% box margins and rules
    boxrule=1pt,
    left=0.2cm,right=0.2cm,top=0.2cm,
    toptitle=0.1cm,
    bottomtitle=0.08cm,
% box fonts
    fonttitle=\bfseries,fontupper=\normalsize
}{rem}

% Proof Box Styles

\title{Group}
\author{Hyoung Joon, Kim}
\date{\today}

\begin{document}

\maketitle
\tableofcontents

\newpage

\section{Group}
\subsection{Introduction and Basic Definitions}
A group is a fundamental algebraic structure consisting of a set equipped with a binary operation that satisfies certain axioms. Groups are widely used in various fields of mathematics, including algebra, geometry, and number theory, as well as in physics and computer science.


Now, we provide the formal definition of a group.
% Definition of Group
\begin{defn}{Group}{GP}
A \emph{group} is a set \( G \) equipped with a binary operation
\(\cdot : G \times G \to G\) satisfying the following axioms:
\begin{enumerate}[label=(\roman*)]
    \item (\emph{Closure}) For all \( a, b \in G \), the result of the operation \( a \cdot b \) is also in \( G \).
    \item (\emph{Associativity}) For all \( a, b, c \in G \), \( (a \cdot b) \cdot c = a \cdot (b \cdot c) \).
    \item (\emph{Identity Element}) There exists an element \( e \in G \) such that for every element \( a \in G \), \( e \cdot a = a \cdot e = a \).
    \item (\emph{Inverse Element}) For each \( a \in G \), there exists an element \( b \in G \) such that \( a \cdot b = b \cdot a = e \), where \( e \) is the identity element.
\end{enumerate}
Furthermore, if the binary operation is commutative, i.e., \(a \cdot b = b \cdot a\) for all \(a, b \in G\), then the group is called an \emph{abelian group}.
\end{defn}

Here, we can weaken the conditions (iii) and (iv) of the \hyperref[defn:GP]{Definition~\ref{defn:GP}} as follows:
% Remark on Weakened Group Definition
\begin{rem}{}{GPWeaken}
Change condition (iii) and (iv) of \hyperref[defn:GP]{Definition~\ref{defn:GP}} to the following:
\begin{enumerate}[label=(\roman*)]
    \item[(iii)$'$] (\emph{Left Identity Element}) There exists an element \( e \in G \) such that for every element \( a \in G \), \( e \cdot a = a \).
    \item[(iv)$'$] (\emph{Left Inverse Element}) For each \( a \in G \), there exists an element \( b \in G \) such that \( b \cdot a = e \), where \( e \) is the identity element.
\end{enumerate}
Then, the set \( G \) with the binary operation \(\cdot\) still forms a group.
\end{rem}

\begin{proof}
It is left as an exercise.
\end{proof}

\begin{defn}{Subgroup}{SubGP}
    Let $(G,\ \cdot)$ be a group. A subset $H \in G$ is called a \emph{subgroup} of $G$ if $H$ satisfies the following conditions:
    \begin{enumerate}[label=(\roman*)]
        \item (\emph{Closure}) For all \( a, b \in H \), \( a \cdot b \in H \).
        \item (\emph{Identity Element}) The identity element \( e \in G \) is in \( H \).
        \item (\emph{Inverse Element}) For each \( a \in H \), \(a^{-1} \in H\).
    \end{enumerate} 
\end{defn}

\end{document}